\documentclass[11pt,twoside,openright]{book}

\usepackage[a4paper,inner=27mm,outer=23mm,top=25mm,bottom=28mm]{geometry}
\usepackage[T1]{fontenc}
\usepackage{lmodern}
\usepackage[provide=*,french]{babel}
\usepackage{microtype}
\usepackage{amsmath,amssymb,mathtools}
\usepackage{booktabs,array,tabularx}
\usepackage{graphicx}
\usepackage{xcolor}
\usepackage[most]{tcolorbox}
\usepackage{enumitem}
\usepackage{fancyhdr}
\usepackage{hyperref}

\definecolor{DiderotRed}{HTML}{7A1F2B}
\definecolor{DiderotBlue}{HTML}{126E82}
\definecolor{DiderotMint}{HTML}{E8F3F0}
\definecolor{DiderotGold}{HTML}{C99A2E}
\hypersetup{
  colorlinks=true,
  linkcolor=DiderotRed,
  urlcolor=DiderotBlue,
  citecolor=DiderotBlue,
  pdftitle={Diderot ML - Comprendre la Machine Learning Specialization},
  pdfauthor={Guillaume Harbonnier},
  pdfsubject={Distribution gaussienne et detection d'anomalies},
  pdfkeywords={machine learning, gaussienne, detection d'anomalies, Diderot ML}
}

\newtcolorbox{intuitionbox}{enhanced,breakable,colback=DiderotMint,colframe=DiderotBlue,
  title=Intuition, fonttitle=\bfseries,boxrule=0.7pt,arc=1mm}
\newtcolorbox{attentionbox}{enhanced,breakable,colback=yellow!7,colframe=DiderotGold!80!black,
  title=Attention, fonttitle=\bfseries,boxrule=0.7pt,arc=1mm}
\newtcolorbox{engineeringbox}{enhanced,breakable,colback=blue!3,colframe=DiderotBlue,
  title=Application d'ingenierie, fonttitle=\bfseries,boxrule=0.7pt,arc=1mm}
\newtcolorbox{revisionbox}{enhanced,breakable,colback=red!3,colframe=DiderotRed,
  title=Fiche de revision, fonttitle=\bfseries,boxrule=0.8pt,arc=1mm}

\pagestyle{fancy}
\fancyhf{}
\fancyhead[LE]{\small\textsc{Diderot ML}}
\fancyhead[RO]{\small\nouppercase{\leftmark}}
\fancyfoot[C]{\thepage}
\renewcommand{\headrulewidth}{0.4pt}

\setcounter{tocdepth}{2}
\setcounter{secnumdepth}{3}
\setlist{itemsep=0.25em,topsep=0.5em}

\begin{document}

\frontmatter
\begin{titlepage}
  \centering
  \vspace*{23mm}
  {\sffamily\bfseries\color{DiderotRed}\fontsize{30}{35}\selectfont
   Diderot ML\par}
  \vspace{5mm}
  {\sffamily\bfseries\fontsize{22}{27}\selectfont
   Comprendre la Machine Learning Specialization\par}
  \vspace{6mm}
  {\Large Des formules a l'ingenierie des systemes\par}
  \vfill
  \begin{tcolorbox}[width=.82\textwidth,colback=DiderotMint,colframe=DiderotBlue,arc=1mm]
  \centering
  \textbf{Volume de travail v0.1.0}\par
  Distribution gaussienne et fondements de la detection d'anomalies
  \end{tcolorbox}
  \vfill
  {\large Guillaume Harbonnier\par}
  \vspace{2mm}
  {\small 19 juillet 2026\par}
\end{titlepage}

\chapter*{Preface}
\addcontentsline{toc}{chapter}{Preface}

Ce livre est un compagnon d'apprentissage personnel. Son objectif n'est pas de
reproduire un cours existant, mais de transformer chaque notion rencontree en
un savoir durable : intuition, lecture des formules, demonstration, experience
numerique, limites et application a l'ingenierie.

Les chapitres sont rediges en francais, avec une attention particuliere aux
questions qui surgissent naturellement pendant l'apprentissage. Pourquoi une
densite n'est-elle pas une probabilite? Pourquoi $\pi$ apparait-il dans une
formule sans cercle visible? Pourquoi la moyenne minimise-t-elle les distances
carrees? Ces questions ne sont pas des digressions. Elles constituent souvent
le chemin le plus direct vers une comprehension profonde.

La version presente inaugure l'ouvrage avec la distribution gaussienne et la
detection d'anomalies. Le depot associe conserve le code, les tests, la
tracabilite des sources privees et le journal editorial. Les contenus tiers ne
sont pas republies; seules des explications originales et des metadonnees de
tracabilite figurent dans l'ouvrage.

\vspace{1em}
\begin{flushright}
Guillaume Harbonnier\\
Diderot ML, 2026
\end{flushright}


\tableofcontents

\mainmatter
\part{Apprentissage non supervise}
\chapter{La gaussienne, brique de la detection d'anomalies}
\label{chap:gaussian}

\section{Le probleme que nous voulons resoudre}

Un systeme produit continuellement des mesures : temps de reponse, temperature,
qualite d'une capture biometrique, nombre de tentatives, score de comparaison ou
age d'un credential. La plupart de ces mesures decrivent le fonctionnement
habituel. Quelques-unes peuvent cependant etre rares, inattendues ou
incompatibles avec le regime connu.

La detection d'anomalies commence par une question simple :

\begin{quote}
Quelle distribution decrit convenablement les observations normales, et a quel
point une nouvelle observation est-elle compatible avec cette distribution?
\end{quote}

Pour une premiere caracteristique numerique $X$, nous utiliserons une
distribution gaussienne, aussi appelee loi normale :
\begin{equation}
X\sim\mathcal N(\mu,\sigma^2).
\end{equation}

Cette notation se lit : ``$X$ suit une loi normale de moyenne mu et de variance
sigma carre.''

\section{Variable aleatoire, moyenne et dispersion}

Une variable aleatoire $X$ associe une valeur numerique a une experience dont
le resultat exact n'est pas connu a l'avance. Les observations disponibles sont
notees
\begin{equation}
x^{(1)},x^{(2)},\ldots,x^{(m)}.
\end{equation}
L'exposant entre parentheses est un indice d'exemple, pas une puissance.

La moyenne $\mu$ localise le centre du comportement :
\begin{equation}
\mu=\mathbb E[X].
\end{equation}
La variance mesure la dispersion quadratique autour du centre :
\begin{equation}
\sigma^2=\mathbb E[(X-\mu)^2].
\end{equation}
L'ecart-type est la racine carree de la variance :
\begin{equation}
\sigma=\sqrt{\sigma^2}.
\end{equation}

\begin{attentionbox}
$\sigma$ est l'ecart-type et $\sigma^2$ est la variance. Si une latence est
mesuree en millisecondes, $\sigma$ s'exprime en millisecondes tandis que
$\sigma^2$ s'exprime en millisecondes carrees.
\end{attentionbox}

\section{La densite gaussienne}

La densite d'une gaussienne univariee est
\begin{equation}
p(x;\mu,\sigma^2)=
\frac{1}{\sqrt{2\pi}\,\sigma}
\exp\left(-\frac{(x-\mu)^2}{2\sigma^2}\right).
\label{eq:gaussian}
\end{equation}

Elle contient deux mecanismes complementaires. Le terme
\begin{equation}
\exp\left(-\frac{(x-\mu)^2}{2\sigma^2}\right)
\end{equation}
fait decroitre la densite lorsque $x$ s'eloigne de $\mu$. Le facteur
\begin{equation}
\frac{1}{\sqrt{2\pi}\sigma}
\end{equation}
normalise la courbe afin que son aire totale soit egale a un :
\begin{equation}
\int_{-\infty}^{+\infty}p(x)\,dx=1.
\end{equation}

\subsection{Densite et probabilite ne sont pas synonymes}

Pour une variable continue, $p(x)$ est une hauteur de densite. La probabilite
d'appartenir a un intervalle est une aire :
\begin{equation}
\Pr(a\le X\le b)=\int_a^b p(x)\,dx.
\end{equation}
Pour une classe etroite de largeur $\Delta x$ centree en $x$,
\begin{equation}
\Pr\left(x-\frac{\Delta x}{2}\le X\le x+\frac{\Delta x}{2}\right)
\approx p(x)\Delta x.
\end{equation}
Avec $m$ observations, l'effectif attendu dans cette classe est donc environ
$m p(x)\Delta x$.

\begin{intuitionbox}
Un histogramme normalise est une approximation empirique de la densite. Quand
le nombre d'observations augmente et que les classes deviennent suffisamment
fines, sa forme se rapproche de la courbe theorique $p(x)$.
\end{intuitionbox}

\section{Effet de la moyenne et de l'ecart-type}

La figure~\ref{fig:gaussian-comparison} compare quatre distributions. La
moyenne deplace la courbe; l'ecart-type change sa largeur.

\begin{figure}[htbp]
\centering
\includegraphics[width=.92\textwidth]{figures/gaussian-comparison.pdf}
\caption{Effet de $\mu$ et $\sigma$ sur la densite gaussienne. Toutes les aires
valent un. Figure originale generee par le code du depot.}
\label{fig:gaussian-comparison}
\end{figure}

Le maximum de la densite se trouve en $x=\mu$ :
\begin{equation}
p(\mu)=\frac{1}{\sqrt{2\pi}\sigma}.
\end{equation}
Ainsi, pour $\sigma=1$, la hauteur maximale vaut environ $0{,}399$. Pour
$\sigma=0{,}5$, elle vaut environ $0{,}798$. Une courbe plus etroite doit etre
plus haute puisque son aire reste egale a un.

\begin{center}
\begin{tabular}{@{}cccc@{}}
\toprule
$\mu$ & $\sigma$ & $\sigma^2$ & Effet principal \\
\midrule
$0$ & $1$ & $1$ & reference centree reduite \\
$0$ & $0{,}5$ & $0{,}25$ & courbe etroite et haute \\
$0$ & $2$ & $4$ & courbe large et basse \\
$3$ & $0{,}5$ & $0{,}25$ & meme forme, centre deplace \\
\bottomrule
\end{tabular}
\end{center}

\section{Pourquoi \texorpdfstring{$\pi$}{pi} apparait-il sans cercle visible?}

Le facteur $\sqrt{2\pi}$ provient de l'integrale gaussienne. Considerons
\begin{equation}
I=\int_{-\infty}^{+\infty}e^{-x^2/2}\,dx.
\end{equation}
Cette integrale resiste aux primitives elementaires. L'idee decisive est de la
mettre au carre :
\begin{align}
I^2
&=\left(\int_{-\infty}^{+\infty}e^{-x^2/2}\,dx\right)
  \left(\int_{-\infty}^{+\infty}e^{-y^2/2}\,dy\right)\\
&=\iint_{\mathbb R^2}e^{-(x^2+y^2)/2}\,dx\,dy.
\end{align}
Le cercle cache apparait maintenant : $x^2+y^2=r^2$. En coordonnees polaires,
$dx\,dy=r\,dr\,d\theta$, donc
\begin{align}
I^2
&=\int_0^{2\pi}\int_0^{+\infty}e^{-r^2/2}r\,dr\,d\theta\\
&=2\pi\left[-e^{-r^2/2}\right]_{0}^{+\infty}=2\pi.
\end{align}
Par consequent,
\begin{equation}
I=\sqrt{2\pi}.
\end{equation}

\begin{intuitionbox}
Le cercle n'est pas visible dans la courbe unidimensionnelle. Il apparait quand
on duplique l'integrale et que la symetrie quadratique $x^2+y^2$ devient une
symetrie circulaire. Le nombre $\pi$ vient du tour complet de l'angle
$0\le\theta\le2\pi$.
\end{intuitionbox}

\section{Gaussienne centree, reduite et score standardise}

Une loi normale est centree lorsque $\mu=0$. Elle est centree reduite lorsque
\begin{equation}
Z\sim\mathcal N(0,1).
\end{equation}
Toute observation gaussienne peut etre standardisee :
\begin{equation}
z=\frac{x-\mu}{\sigma}.
\end{equation}
Le score $z$ indique a combien d'ecarts-types l'observation se trouve du centre.
La densite s'ecrit alors
\begin{equation}
p(x)=\frac{1}{\sigma\sqrt{2\pi}}e^{-z^2/2}.
\end{equation}
C'est donc la distance quadratique normalisee $z^2$ qui controle la chute de la
densite.

\section{Estimer les parametres a partir des donnees}

Dans une application, $\mu$ et $\sigma^2$ sont inconnus. A partir des $m$
observations d'entrainement, nous utilisons
\begin{equation}
\widehat\mu=\frac{1}{m}\sum_{i=1}^{m}x^{(i)},
\label{eq:mu-hat}
\end{equation}
et
\begin{equation}
\widehat\sigma^2=
\frac{1}{m}\sum_{i=1}^{m}\left(x^{(i)}-\widehat\mu\right)^2.
\label{eq:sigma-hat}
\end{equation}
Le chapeau rappelle qu'il s'agit d'estimations calculees sur un echantillon.

\subsection{Exemple numerique}

Pour les valeurs $8,9,10,11,12$,
\begin{equation}
\widehat\mu=10.
\end{equation}
Les ecarts au centre sont $-2,-1,0,1,2$, et leurs carres sont $4,1,0,1,4$.
La variance estimee est
\begin{equation}
\widehat\sigma^2=\frac{4+1+0+1+4}{5}=2,
\end{equation}
d'ou $\widehat\sigma=\sqrt{2}\approx1{,}414$. Le modele ajuste est
$\mathcal N(10,2)$, le second parametre etant la variance.

\section{Maximum de vraisemblance}

Les formules~\eqref{eq:mu-hat} et~\eqref{eq:sigma-hat} sont les estimateurs du
maximum de vraisemblance. On suppose les observations independantes et issues
d'une meme gaussienne. Leur vraisemblance conjointe est
\begin{equation}
L(\mu,\sigma^2)=\prod_{i=1}^{m}p(x^{(i)};\mu,\sigma^2).
\end{equation}
Nous cherchons les parametres rendant les donnees observees aussi vraisemblables
que possible :
\begin{equation}
(\widehat\mu,\widehat\sigma^2)=
\arg\max_{\mu,\sigma^2}L(\mu,\sigma^2).
\end{equation}

Le logarithme transforme le produit en somme :
\begin{equation}
\log L=-\frac{m}{2}\log(2\pi)-m\log\sigma
-\frac{1}{2\sigma^2}\sum_{i=1}^{m}(x^{(i)}-\mu)^2.
\end{equation}
Annuler les derivees par rapport a $\mu$ et $\sigma^2$ redonne les deux
estimateurs precedents.

\subsection{Le pont avec K-means}

Maximiser la vraisemblance par rapport a $\mu$ revient notamment a minimiser
\begin{equation}
\sum_{i=1}^{m}(x^{(i)}-\mu)^2.
\end{equation}
Le minimum est atteint par la moyenne. C'est exactement la raison pour laquelle
K-means place un centroide a la moyenne des points qui lui sont affectes. Le
modele gaussien ajoute au centre une mesure de dispersion et une interpretation
probabiliste.

\subsection{Pourquoi parfois \texorpdfstring{$m-1$}{m-1}?}

L'estimateur de maximum de vraisemblance divise par $m$. En statistique
descriptive, on rencontre aussi
\begin{equation}
s^2=\frac{1}{m-1}\sum_{i=1}^{m}(x^{(i)}-\bar x)^2.
\end{equation}
La division par $m-1$ corrige le biais d'estimation de la variance d'une
population lorsque la moyenne a ete estimee sur le meme echantillon. Pour de
grands echantillons, la difference pratique devient faible.

\section{De la densite a la decision d'anomalie}

Une fois $\widehat\mu$ et $\widehat\sigma^2$ appris, une nouvelle observation
$x_{\mathrm{new}}$ recoit la densite
\begin{equation}
p(x_{\mathrm{new}};\widehat\mu,\widehat\sigma^2).
\end{equation}
La regle de decision utilise un seuil $\varepsilon$ :
\begin{equation}
p(x_{\mathrm{new}})<\varepsilon
\quad\Longrightarrow\quad
\text{anomalie potentielle}.
\end{equation}

Une faible densite ne prouve ni fraude ni panne. Elle signifie seulement que
l'observation appartient a une region peu compatible avec le modele normal.
Une analyse de contexte reste necessaire.

\begin{engineeringbox}
Pour un systeme d'identite, une anomalie peut provenir d'une fraude, d'un capteur
defaillant, d'une integration cliente incorrecte, d'un changement de population
ou d'un utilisateur legitime confronte a un processus inadapte. Un moteur de
risque responsable doit donc distinguer detection statistique, qualification de
cause et decision operationnelle, notamment pour eviter l'exclusion.
\end{engineeringbox}

\section{Passage a plusieurs caracteristiques}

Une application reelle observe un vecteur
\begin{equation}
\mathbf x=(x_1,x_2,\ldots,x_n).
\end{equation}
Une premiere extension ajuste une gaussienne par caracteristique :
\begin{equation}
x_j\sim\mathcal N(\mu_j,\sigma_j^2),
\end{equation}
puis combine les densites :
\begin{equation}
p(\mathbf x)=\prod_{j=1}^{n}p(x_j;\mu_j,\sigma_j^2).
\end{equation}
Cette construction sera developpee dans le chapitre suivant. Elle repose sur
une hypothese d'independance ou, plus prudemment, sur l'idee que les dependances
non modelisees restent acceptables pour la decision recherchee.

En pratique, le calcul numerique utilise souvent la log-densite :
\begin{equation}
\log p(\mathbf x)=\sum_{j=1}^{n}\log p(x_j),
\end{equation}
afin d'eviter qu'un produit de nombreux petits nombres ne soit arrondi a zero.

\section{Limites et attention epistemologique}

\begin{attentionbox}
Une valeur rare n'est pas necessairement une erreur, et une fraude peut etre
statistiquement ordinaire. Le resultat depend des variables choisies, de la
representativite des donnees dites normales, de la stabilite du regime et du
seuil de decision. Une contamination du jeu d'entrainement par des anomalies
peut deplacer la moyenne, gonfler la variance et rendre le detecteur moins
sensible.
\end{attentionbox}

Il faut notamment verifier :
\begin{itemize}
  \item la forme empirique de chaque distribution;
  \item la presence de valeurs extremes dans l'entrainement;
  \item les correlations entre caracteristiques;
  \item les differences entre capteurs, populations et periodes;
  \item la derive temporelle;
  \item le cout respectif des faux positifs et des faux negatifs.
\end{itemize}

\section{Exercices}

\subsection*{Exercice 1 - Lecture de notation}
Lire a voix haute $X\sim\mathcal N(3,0{,}25)$ et donner $\mu$, $\sigma$ et
$\sigma^2$.

\subsection*{Exercice 2 - Estimation}
Pour $x=(2,4,4,5,5)$, calculer $\widehat\mu$, $\widehat\sigma^2$ avec la
division par $m$, puis $\widehat\sigma$.

\subsection*{Exercice 3 - Effet de l'echelle}
Une latence suit $\mathcal N(100,10^2)$. Calculer le score $z$ associe a
$x=130$ ms et expliquer son interpretation sans conclure automatiquement a une
panne.

\subsection*{Exercice 4 - Cercle cache}
Expliquer en quatre lignes pourquoi $\pi$ apparait dans la constante de
normalisation de la gaussienne.

\subsection*{Exercice 5 - Conception d'un detecteur}
Proposer cinq caracteristiques pour detecter des comportements inhabituels dans
un ecosysteme d'identite. Pour chacune, identifier au moins un cas rare mais
legitime.

\section{Reponses breves}

\begin{enumerate}
  \item Moyenne $3$, variance $0{,}25$, ecart-type $0{,}5$.
  \item $\widehat\mu=4$; ecarts carres $4,0,0,1,1$;
  $\widehat\sigma^2=6/5=1{,}2$ et $\widehat\sigma\approx1{,}095$.
  \item $z=(130-100)/10=3$. La mesure est eloignee de trois ecarts-types du
  centre, donc inhabituelle sous ce modele, mais sa cause reste a qualifier.
  \item Le carre de l'integrale devient une integrale en deux dimensions. Le
  terme $x^2+y^2$ possede une symetrie circulaire. Les coordonnees polaires
  introduisent un angle parcourant $0$ a $2\pi$.
  \item La reponse doit distinguer signal statistique, contexte et cause.
\end{enumerate}

\begin{revisionbox}
\begin{align*}
X&\sim\mathcal N(\mu,\sigma^2),\\
p(x)&=\frac{1}{\sqrt{2\pi}\sigma}
e^{-(x-\mu)^2/(2\sigma^2)},\\
\widehat\mu&=\frac1m\sum_i x^{(i)},\\
\widehat\sigma^2&=\frac1m\sum_i(x^{(i)}-\widehat\mu)^2,\\
z&=\frac{x-\mu}{\sigma},\\
p(x)&<\varepsilon\quad\Rightarrow\quad\text{anomalie potentielle}.
\end{align*}
$\mu$ deplace la courbe; $\sigma$ change sa largeur; l'aire totale vaut un;
$p(x)$ est une densite; une anomalie statistique n'est pas une cause.
\end{revisionbox}


\part{Carnet de progression}
\chapter{Carte de l'ouvrage et methode de construction}

Le livre suivra les trois grands ensembles de la specialisation : apprentissage
supervise, algorithmes avances, puis apprentissage non supervise, recommandation
et apprentissage par renforcement. Chaque notion sera reliee aux fils
transverses utiles a l'ingenierie : biometrie, identite, qualite, risque,
observabilite et apprentissage continu.

\section{Un chapitre en sept couches}

Chaque chapitre vise sept niveaux de comprehension :
\begin{enumerate}
  \item le probleme concret;
  \item l'intuition visuelle;
  \item la notation et sa lecture orale;
  \item la derivation mathematique;
  \item l'implementation reproductible;
  \item les limites et les risques d'interpretation;
  \item le transfert vers les projets d'ingenierie.
\end{enumerate}

\section{Prochaines integrations}

Les prochaines versions consolideront d'abord les notes deja disponibles sur
K-means, les arbres de decision et les ensembles. Les nouveaux chapitres seront
ensuite ajoutes au fil de la progression, avec une fiche de revision et un TP
pour chaque bloc conceptuel important.



\backmatter
\chapter{Notation minimale}

\begin{tabularx}{\textwidth}{@{}lX@{}}
\toprule
Symbole & Signification \\
\midrule
$m$ & nombre d'exemples d'entrainement \\
$n$ & nombre de caracteristiques \\
$x^{(i)}$ & valeur du $i$-eme exemple; ce n'est pas une puissance \\
$x_j^{(i)}$ & caracteristique $j$ de l'exemple $i$ \\
$\mu$ & moyenne de la population ou parametre de centre \\
$\widehat\mu$ & moyenne estimee sur l'echantillon \\
$\sigma$ & ecart-type \\
$\sigma^2$ & variance \\
$p(x)$ & densite au point $x$ \\
$\varepsilon$ & seuil de decision d'anomalie \\
$\mathcal N(\mu,\sigma^2)$ & loi normale de moyenne $\mu$ et variance $\sigma^2$ \\
\bottomrule
\end{tabularx}



\end{document}
